\subsubsection{NULL Handling}

\texttt{NULL} means unknown, so it propagates through arithmetic and poisons
comparisons. The two functions below are the standard guards.

\paragraph{COALESCE and NULLIF.} \texttt{COALESCE(a, b, c, ...)} returns the
first non-\texttt{NULL} argument, evaluated left to right, or \texttt{NULL} if
all are \texttt{NULL}. The common use is a fallback that keeps a \texttt{NULL}
from poisoning arithmetic: \texttt{salary + commission} is \texttt{NULL}
whenever \texttt{commission} is, so write \texttt{salary + COALESCE(commission, 0)}.
It also supplies a default down a priority list, as in
\texttt{COALESCE(nickname, first\_name, 'Unknown')}. It is ANSI standard,
unlike the vendor-specific \texttt{IFNULL} (MySQL) and \texttt{ISNULL}
(SQL Server), and it accepts any number of arguments.

\texttt{NULLIF(a, b)} is the mirror image: it returns \texttt{NULL} when
\texttt{a = b}, otherwise \texttt{a}. Its standard use guards division,
\texttt{x / NULLIF(y, 0)} yields \texttt{NULL} instead of raising a
divide-by-zero when \texttt{y} is 0.

\begin{lstlisting}[style=sql, caption={NULLIF guards division, COALESCE supplies defaults}]
SELECT
    order_id,
    revenue / NULLIF(quantity, 0)  AS unit_price,  -- NULL, not an error, when quantity = 0
    COALESCE(discount, 0)          AS discount      -- NULL discount treated as 0
FROM orders;
\end{lstlisting}

\vspace{-1ex}

A subtlety with aggregates: \texttt{SUM}, \texttt{AVG}, and
\texttt{COUNT(col)} skip \texttt{NULL}s. So \texttt{AVG(commission)} averages
only the non-\texttt{NULL} rows, while \texttt{AVG(COALESCE(commission, 0))}
counts the \texttt{NULL}s as zeros and enlarges the denominator. 

\paragraph{Gotchas.}
\begin{itemize}
  \item \textbf{\texttt{NOT IN} with NULLs:} if the subquery returns any
  \texttt{NULL}, \texttt{x NOT IN (...)} is never true and the query returns
  zero rows. \texttt{x NOT IN (1, NULL)} expands to
  \texttt{NOT (x=1 OR x=NULL)}, and the \texttt{NULL} comparison poisons the
  result to unknown. Use \texttt{NOT EXISTS}, which is null-safe.
  \item \textbf{Integer division truncates:} \texttt{5 / 2} yields \texttt{2}
  because the fractional part is discarded. Cast one operand:
  \texttt{5::numeric / 2} or \texttt{5.0 / 2}. Bites hardest in ratio and
  percentage columns.
  \item \textbf{\texttt{COUNT(col)} vs \texttt{COUNT(*)}:} the former skips
  \texttt{NULL}s in \texttt{col}, the latter counts every row. A common source
  of off-by-$N$ errors.
\end{itemize}
